%% ----------------------------------------------------------------------------
%% ReScience X article template
%% Copyright (c) 2022 N.P. Rougier, E.B. Roesch
%% Released under a Creative Commons Attribution 4.0 International license.
%% ----------------------------------------------------------------------------

%% --- Top left header --------------------------------------------------------
\begin{tikzpicture}[remember picture, overlay]
  \node [shape=rectangle, draw=darkred, fill=darkred, yshift=-34mm,
        anchor=north west, minimum width=3.75cm, minimum height=10mm]
        at (current page.north west) {};
  \node [text=white, anchor=center, yshift=-39mm, xshift=1.875cm]
         at (current page.north west) {\small \SpaceRoboto RESCIENCE X};
\end{tikzpicture}


%% --- Title -------------------------------------------------------------------
{\let\newpage\relax\maketitle} \maketitle


%% --- Margin information -----------------------------------------------------
\marginnote{
  \footnotesize \sffamily
  \textbf{Edited by}\\
  \ifdefempty{\editorNAME}{\textcolor{darkgray}{(Editor)}}
             {\editorNAME\ifdefempty{\editorORCID}{}{$^{\orcid{\editorORCID}}$}}\\
  ~\\
  %%
  \textbf{Reviewed by}\\
  \ifdefempty{\reviewerINAME}{\textcolor{darkgray}{(Reviewer 1)}}
             {\reviewerINAME\ifdefempty{\reviewerIORCID}{}{$^{\orcid{\editorORCID}}$}}\\
  \ifdefempty{\reviewerIINAME}{\textcolor{darkgray}{(Reviewer 2)}}
             {\reviewerIINAME\ifdefempty{\reviewerIIORCID}{}{$^{\orcid{\editorORCID}}$}}\\
  ~\\
  %%           
  \textbf{Received}\\
  \ifdefempty{\dateRECEIVED}{---}{\dateRECEIVED}\\
  ~\\
  %%
  \textbf{Published}\\
  \ifdefempty{\datePUBLISHED}{---}{\datePUBLISHED}\\
  ~\\
  %%
  \textbf{DOI}\\
  \ifdefempty{\articleDOI}{---}{\articleDOI}
}
 

%% --- Bottom statement -------------------------------------------------------
\newcommand{\container}[1]{\def\@container{#1}}
\begin{container}
  \afterpage {
    \begin{statement}
      \scriptsize \sffamily
      \hrule \vskip .5em
      Copyright © \articleYEAR~\authorsABBRV,
      released under a Creative Commons Attribution 4.0 International license.
      
      Correspondence should be addressed to
      \contactNAME~(\href{mailto:\contactEMAIL}{\contactEMAIL})
  
      The authors have declared that no competing interests exists.

      \ifdefempty{\codeURL}{}
      {Code is available at
      \href{\codeURL}{\codeURL}\ifdefempty{\codeDOI}{.}{ -- DOI \doi{\codeDOI}.}}

      \ifdefempty{\dataURL}{}
      {Data is available at
      \href{\dataURL}{\dataURL}\ifdefempty{\dataDOI}{.}{ -- DOI \doi{\dataDOI}.}}

      \ifdefempty{\reviewURL}{}
     {Open peer review is available at \href{\reviewURL}{\reviewURL}.}
    \end{statement}
  }
\end{container}


%% --- Abstract ---------------------------------------------------------------
%% \ifdefempty{\articleABSTRACT}{}{
%%   {\small \sffamily \textbf{Abstract} \articleABSTRACT \par}
%% }


%% --- Chymeric section
\setlength\parskip{0.5em plus 0.1em minus 0.2em}


%% --- Replication ------------------------------------------------------------
\ifdefempty{\replicationBIB}{}{
  {\small \sffamily \bfseries A replication of}
  \ifdefempty{\replicationURL}
       {\href{http://doi.org/\replicationDOI}{\small \sffamily \replicationBIB}}
       {\href{\replicationURL}{\small \sffamily \replicationBIB}}.
\par }


\title{
	\href{https://github.com/con/opfvta-reexecution}{
		\Large github.com/con/opfvta-reexecution
	}\\\vspace{.15em}
	Neuroimaging Article Reexecution and Reproduction Assesment System
}
\author{
	Horea-Ioan Ioanas$^{1}$,
	Austin Macdonald$^{1}$,
	Yaroslav O. Halchenko$^{1}$
}
\institute[CON]{$^{1}$Center for Open Neuroscience, Department of Psychological and Brain Sciences, Dartmouth College}
\date{\today}

\newlength{\columnheight}
\setlength{\columnheight}{0.881\textheight}

\begin{document}

\begin{frame}
\begin{columns}
	\begin{column}{.42\textwidth}
		\begin{beamercolorbox}[center]{postercolumn}
			\begin{minipage}{.98\textwidth}  % tweaks the width, makes a new \textwidth
				\parbox[t][\columnheight]{\textwidth}{ % must be some better way to set the the height, width and textwidth simultaneously
					\begin{myblock}{Abstract}
						\section{Abstract}

As complexity increases, the value of research articles is inextricably contingent on data analysis results which substantiate their claims.
Unlike data production steps, data analysis steps lend themselves to a higher standard of both transparency and repeated operator-free execution.
This higher standard can be approached via fully reexecutable research outputs, which contain the entire instruction set for end-to-end generation of an entire article solely from the earliest feasible provenance point, in a programatically executable format.
% TODO; maybe put “as a starting point” instead of “in order”
In this study, we make use of a peer-reviewed neuroimaging article, which provides complete but fragile reexecution instructions, in order to formulate a new reexecution model which is both robust and portable.
We render this model modular as a core design aspect, so that reexecutable article code, data, and environment specifications can be easily substituted or adapted.
In conjunction with this model, which forms the main reusable product of this study, we detail the core challenges with full article reexecution and specify a number of best practices which permitted us to mitigate them.
We further show how these capabilities can subsequently be used to provide reproducibility assesments, both via simple statistical metrics and by visually highlighting divergent elements for human inspection.
We argue that reexecutable articles are thus a feasible best practice, the usage of which can greatly enhance the understanding of data analysis variability.
Lastly, we comment at length on the outlook for reexecutable resource and encourage re-use and derivation of the model produced herein.

% Maybe also mention, though ideally somewhere else:
% * The comparison article we thus produce, i.e. this article, is similarly reexecutable.
% * ... a collection of best practices, including YODA principles for data management, containers for preservation, Gentoo for long term flexibility.
% * This is a very good point, though again, maybe not for the abstract “The inability to re-use code due to instability thus limits the lasting value of the work as a repository of procedural knowledge”.
% *We document a number of prominent difficulties with de novo article generation, arising from the rapid evolution of extrinsic tools, and from nondeterministic data analysis procedures.


					\end{myblock}\vfill
					\vspace{-0.3em}
					\begin{myblock}{Workflow}
						\vspace{0.5em}
						\begin{figure}
							\includegraphics[width=0.99\textwidth]{figs/workflow.pdf}
							\caption{
								The reexecution system encompasses both the original article (first target), and the “meta-article” publishing materials (article manuscript, as well as this poster), the latter of which takes user- and developer-submitted reexecution results as an input for the reproduction quality assessment.
							}
						\end{figure}
					\end{myblock}\vfill
					\vspace{-0.3em}
					\begin{myblock}{Topoplogy}
						\vspace{0.5em}
						\begin{figure}
							\captionsetup{width=.9\linewidth}
							\includegraphics[width=0.99\textwidth]{figs/topology.pdf}
							\caption{
								The reexecution workflow is supported by a resource topology in which reexecution code (first box), “meta-article” code (second box), reexecution resources (third box), and the reexecution output record (last box) are separated at the top directory level.
								The figure depicts direcotry trees via nested boxes, with external resources automatically fetched as via the reexecution code being highlighted in orange.
								The green highlighted article represents a sample reexecution output, and the blue hignlighted article represents the manuscript, an analogous output to this poster generated in the same directory.
							}
							\label{fig:workflow}
						\end{figure}
					\end{myblock}\vfill
					\begin{myblock}{Best Practice Guidelines}
						\vspace{0.5em}
						As part of setting up an encompassing reexecution system, we formulate a number of best practices, including:
						\begin{itemize}
							\item \textbf{Errors should be fatal more often than not.}\\
								\colorbox{elg}{\texttt{set -eu}}, prepended to POSIX shell scripts, will ensure that workflows fail when a subcommand does, or when an encountered variable is undefined.
							%\item \textbf{Avoid assuming or hard-coding absolute paths to resources.}\\
							\item \textbf{Avoid assuming a directory context for execution.}\\
								\colorbox{elg}{\texttt{cd \textquotedbl\$(dirname \textquotedbl\$0\textquotedbl)\textquotedbl}}, prepended to POSIX shell scripts, will ensure that in complex workflows scripts can operate relative to their location directory context and not the execution context.
							\item \textbf{Workflow granularity greatly benefits efficiency.}\\
								While the underlying execution system of the target article, RepSeP \cite{repsep} separates data analysis into two distinct (voxel-space “low iteration” and top-level “high iteration”) steps, further granularity can benefit debugging, particularly in container environments.
							%\item Container image size should be kept small.
							\item \textbf{Resources should be bundled into a DataLad superdataset.}\\
								Resource bundling, with usage of submodules for external resources (as seen in \cref{fig:workflow}) allows management of required resources via Git and associated technologies, such as DataLad \cite{datalad} — this is known as the YODA principle \cite{yoda}.
							%\item Containers should fit the scope of the underlying workflow steps.
							%\item Do not write debug-relevant data inside the container
							%\item Parameterize scratch directories
							\item \textbf{Dependency versions inside container environments should be frozen as soon as feasible.}\\
								This is best accomplished via a package manager which uses version tracking for its software provision index; in Gentoo Linux, used here on account of broad provision of neuroscience packages \cite{ng}, this can be done via: \colorbox{elg}{\texttt{cd /.../myrepo; git fetch origin \$myhash; git checkout \$myhash}}.
						\end{itemize}
					\end{myblock}\vfill
					\begin{myblock}{References}
						\vspace{0.5em}
						\begin{minipage}{.3\textwidth}
							\begin{figure}
								\includegraphics[width=0.9\textwidth]{figs/qr.eps}
							\end{figure}
						\end{minipage}
						\begin{minipage}{.69\textwidth}
							\scriptsize
                                                        \bibliographystyle{ieeetr}
                                                        \bibliography{bibliography}
						\end{minipage}
					\end{myblock}\vfill
		}\end{minipage}\end{beamercolorbox}
	\end{column}
	\begin{column}{.59\textwidth}
		\begin{beamercolorbox}[center]{postercolumn}
			\begin{minipage}{.98\textwidth} % tweaks the width, makes a new \textwidth
				\parbox[t][\columnheight]{\textwidth}{ % must be some better way to set the the height, width and textwidth simultaneously
					\begin{myblock}{Reproduction Assessment Showcase}
						\vspace{-0.45em}
						\begin{minipage}{.58\textwidth}
						\begin{figure}
							\includegraphics[width=0.95\textwidth]{figs/diff_pages.pdf}
							\vspace{0.2em}
							\caption{
								Page-wise pixel difference comparison across multiple reexecutions in different environments indicates consistency of variability in both extent and location.
							}
							\label{fig:ras_s}
						\end{figure}
						\begin{figure}
							\fbox{\includegraphics[width=0.9\textwidth]{figs/diff_fig.pdf}}
							\vspace{0.2em}
							\caption{
								One notable source of variability are data plots, where it can be observed that even as data points vary to an almost full extent, statistical summaries can remain constant.
							}
						\end{figure}
						\end{minipage}
						\hfill
						\begin{minipage}{.38\textwidth}
						\vspace{1.3em}
						\begin{figure}
							\fbox{\includegraphics[width=0.9\textwidth]{figs/diff_text.pdf}}
							\vspace{0.5em}
							\caption{
								Text differences in statistical summaries account for a small proportion of pixel differences, but can remain well-localized instead of spreading via test shift if statistical summaries are appropriately trimmed down to a constant length.
							}
						\end{figure}
						\begin{figure}
							\fbox{\includegraphics[width=0.9\textwidth]{figs/diff_date.pdf}}
							\vspace{0.5em}
							\caption{
								A good litmus test for monitoring differences (accounting for the baseline difference in \cref{fig:ras_s}) is the datestamp of the reexecution, which should always be expected to differ from the manuscript.
							}
						\end{figure}
						\end{minipage}
					\end{myblock}\vfill
					\begin{myblock}{Full Document Comparison}
						\vspace{0.25em}
						Reproduction assessment is based on \textit{full} document “diffs”.
						The following figures are excerpts, with tinted highlighting (blue for the original manuscript, and orange for reexecution result).
						First row pages exemplify inline statistical differences and second row pages exemplify figure differences.
						Differing sections are highlighted with a red left-hand marking.\\
						\vspace{0.75em}
						\fbox{\includegraphics[page=4,scale=1.04]{data/marked_paperdiff_singularity_20230908122618.pdf}}
						\fbox{\includegraphics[page=5,scale=1.04]{data/marked_paperdiff_singularity_20230908122618.pdf}}
						\\
						\vspace{.3em}
						\fbox{\includegraphics[page=13,scale=1.04]{data/marked_paperdiff_singularity_20230908122618.pdf}}
						\fbox{\includegraphics[page=14,scale=1.04]{data/marked_paperdiff_singularity_20230908122618.pdf}}
					\end{myblock}\vfill
		}\end{minipage}\end{beamercolorbox}
	\end{column}
\end{columns}
\end{frame}
\end{document}
